% ============================================================
% §2.6 · Detalles de Implementación — Manual Técnico K-QGMIP
% ============================================================

\section{Detalles de Implementación}

% ─────────────────────────────────────────────────────────────────────────────
\subsection{Métodos Principales}
% ─────────────────────────────────────────────────────────────────────────────

A continuación se documentan los métodos públicos y funciones auxiliares más
relevantes del proyecto, con sus firmas exactas, parámetros y valores de retorno.

% ··················································
\subsubsection*{\texttt{KGeoMIP.aplicar\_estrategia}}
% ··················································

\begin{lstlisting}
def aplicar_estrategia(
    self,
    condicion: str,
    alcance: str,
    mecanismo: str,
    tpm: np.ndarray,
    k: int = 2,
    variante: str = "E4",
    estrategia_corte: str = "auto",
) -> Solution:
\end{lstlisting}

\textbf{Descripción:} Halla la k-partición de mínima información integrada
usando la estrategia KGeoMIP con anclaje en GeoMIP. Para $k=1$ retorna
$\varphi=0$ directamente; para $k=2$ retorna la solución de GeoMIP sin cálculo
adicional.

\begin{table}[H]
\centering
\begin{tabular}{llp{5.5cm}}
  \toprule
  \textbf{Parámetro} & \textbf{Tipo} & \textbf{Descripción} \\
  \midrule
  \texttt{condicion}         & \texttt{str}        & Máscara binaria; bit=1 $\to$ nodo candidato \\
  \texttt{alcance}           & \texttt{str}        & Máscara binaria del alcance del subsistema \\
  \texttt{mecanismo}         & \texttt{str}        & Máscara binaria del mecanismo del subsistema \\
  \texttt{tpm}               & \texttt{np.ndarray} & TPM precargada de forma $(2^n, n)$ \\
  \texttt{k}                 & \texttt{int}        & Número de partes ($1 \le k \le 5$) \\
  \texttt{variante}          & \texttt{str}        & \texttt{"E4"} (divisiva guiada) o \texttt{"A"} (aglomerativa) \\
  \texttt{estrategia\_corte} & \texttt{str}        & \texttt{"auto"} / \texttt{"exhaustivo"} / \texttt{"guiado\_S"} \\
  \midrule
  \textbf{Retorno}           & \texttt{Solution}   & \texttt{perdida}, \texttt{particion}, \texttt{tiempo\_total} \\
  \bottomrule
\end{tabular}
\end{table}

\textbf{Comportamiento especial:}
\begin{itemize}
  \item Si la clave \texttt{(condicion, alcance, mecanismo)} coincide con la
        llamada anterior, reutiliza la solución GeoMIP anclada, la matriz $S$
        y todos los cachés, evitando recomputación al barrer
        $k \in \{2,3,4,5\}$ sobre el mismo subsistema.
  \item \texttt{estrategia\_corte = "auto"} usa modo exhaustivo si
        $2^{m-1}-1 \le 4096$ (bloques de hasta 13 nodos), y
        \texttt{guiado\_S} en bloques más grandes.
\end{itemize}

% ··················································
\subsubsection*{\texttt{KQNodes.aplicar\_estrategia}}
% ··················································

\begin{lstlisting}
def aplicar_estrategia(
    self,
    estado_inicial: str,
    condicion: str,
    alcance: str,
    mecanismo: str,
    k: int = 2,
    criterio: str = "C4",
) -> Solution:
\end{lstlisting}

\textbf{Descripción:} Halla la k-partición de mínima información integrada
mediante refinamiento iterativo greedy con el oracle de Queyranne. Para $k=2$
replica exactamente \texttt{QNodes.winner()}.

\begin{table}[H]
\centering
\begin{tabular}{llp{5.5cm}}
  \toprule
  \textbf{Parámetro} & \textbf{Tipo} & \textbf{Descripción} \\
  \midrule
  \texttt{estado\_inicial} & \texttt{str}      & Estado inicial del sistema (string de bits) \\
  \texttt{condicion}       & \texttt{str}      & Máscara binaria de nodos candidatos \\
  \texttt{alcance}         & \texttt{str}      & Máscara binaria del alcance \\
  \texttt{mecanismo}       & \texttt{str}      & Máscara binaria del mecanismo \\
  \texttt{k}               & \texttt{int}      & Número de partes ($2 \le k \le 5$) \\
  \texttt{criterio}        & \texttt{str}      & \texttt{"C4"} ($\min\varphi_{\text{local}}$) o \texttt{"C1"} (máx tamaño) \\
  \midrule
  \textbf{Retorno}         & \texttt{Solution} & \texttt{perdida}, \texttt{particion}, \texttt{tiempo\_total} \\
  \bottomrule
\end{tabular}
\end{table}

% ··················································
\subsubsection*{\texttt{SIA.sia\_preparar\_subsistema}}
% ··················································

\begin{lstlisting}
# Variante GeoMIP (sin estado_inicial; lo gestiona Manager):
def sia_preparar_subsistema(
    self, condicion: str, alcance: str, mecanismo: str
) -> None:

# Variante QNodes:
def sia_preparar_subsistema(
    self, estado_inicial: str, condicion: str,
    alcance: str, mecanismo: str
) -> None:
\end{lstlisting}

\textbf{Descripción:} Método de la clase base \texttt{SIA} (patrón
\textit{Template Method}). Construye el objeto \texttt{System} condicionado y
sustraído desde las máscaras binarias. El subsistema resultante se almacena en
\texttt{sia\_subsistema}; las distribuciones marginales en
\texttt{sia\_dists\_marginales}. No tiene valor de retorno.

% ··················································
\subsubsection*{\texttt{NCube.marginalizar}}
% ··················································

\begin{lstlisting}
def marginalizar(self, ejes: NDArray[np.int8]) -> "NCube":
\end{lstlisting}

\textbf{Descripción:} Colapsa el hipercubo sobre los ejes indicados promediando
las dimensiones correspondientes. Retorna un nuevo \texttt{NCube} con
\texttt{dims} reducidos. Es la operación primitiva de la preparación del
subsistema y de \texttt{\_calcular\_phi\_total}.

\begin{itemize}
  \item \textbf{Parámetro}: \texttt{ejes} — array de índices globales de dimensión
        a eliminar.
  \item \textbf{Retorno}: nuevo \texttt{NCube} con \texttt{data} de forma reducida
        y \texttt{dims} actualizado.
\end{itemize}

% ··················································
\subsubsection*{\texttt{Manager.cargar\_red}}
% ··················································

\begin{lstlisting}
def cargar_red(self) -> np.ndarray:
\end{lstlisting}

\textbf{Descripción:} Lee el CSV de TPM desde \texttt{data/samples/N\{n\}\{muestra\}.csv}
usando \texttt{np.genfromtxt} con delimitador \texttt{";"}.

\begin{itemize}
  \item \textbf{Retorno}: \texttt{ndarray} de forma $(2^n, n)$ con la
        matriz de probabilidad de transición.
  \item \textbf{Excepción}: \texttt{FileNotFoundError} si el archivo CSV
        no existe — propagada sin atrapar para fallo rápido visible.
        La variable de entorno \texttt{GEOMIP\_SAMPLES\_DIR} debe estar
        configurada antes de instanciar \texttt{Manager} (lo hace
        \texttt{main\_kgeomip.py} automáticamente).
\end{itemize}

% ··················································
\subsubsection*{\texttt{\_calcular\_phi\_total}}
% ··················································

\begin{lstlisting}
def _calcular_phi_total(
    particion: list[frozenset[int]],
    sistema: System,
) -> float:
\end{lstlisting}

\textbf{Descripción:} Calcula $\Phi^*$ con una única llamada a
\texttt{emd\_efecto}:
\[
\Phi^* = \mathrm{EMD}\!\left(p(s_{t+1}),\; \bigotimes_{P_i} p_{P_i}\right)
\]
Para cada nodo $d$ en la parte $P_i$, reconstruye su distribución marginal
via \texttt{NCube.marginalizar} eliminando las dimensiones fuera de $P_i$.
Nodos no cubiertos (cuando $N > D_{\text{part}}$) se tratan como singletons.

\textbf{Retorno}: $\Phi^*$ como \texttt{float}. Esta función es la
\textbf{autoridad final} de $\Phi$ — el $\Delta\Phi$ incremental puede diferir
hasta $10^{-6}$ por acumulación de redondeo flotante.

% ··················································
\subsubsection*{\texttt{\_oracle\_restringido}}
% ··················································

\begin{lstlisting}
def _oracle_restringido(
    bloque: frozenset[int],
    N: int,
    D_global: int,
    data_nd: np.ndarray,
    pivot_idx: tuple[int, ...],
    pivot_vals: np.ndarray,
) -> tuple[Callable[[int], float], int]:
\end{lstlisting}

\textbf{Descripción:} Construye el oracle local para el algoritmo de Queyranne
restringido al bloque \texttt{bloque}. Retorna la función
$f_{\text{local}}(\text{mask})$ y la máscara completa del bloque.

El caché de medias \texttt{\_means\_cache} es \textbf{local a cada llamada}
y se reinicia por bloque. Esto es correcto porque dos bloques distintos
$P_i$ y $P_j$ pueden tener la misma máscara numérica local (e.g., \texttt{0b01})
pero con semántica global diferente (decisión D3-02).

% ─────────────────────────────────────────────────────────────────────────────
\subsection{Dependencias Externas}
% ─────────────────────────────────────────────────────────────────────────────

Las dependencias se declaran en \texttt{code/pyproject.toml} bajo
\texttt{[project.dependencies]}. El entorno requiere \textbf{Python $\ge$ 3.12}.

\begin{table}[H]
\centering
\begin{tabular}{llp{6.2cm}}
  \toprule
  \textbf{Biblioteca} & \textbf{Versión mín.} & \textbf{Uso en el proyecto} \\
  \midrule
  \texttt{numpy}        & 2.0.2  & Álgebra lineal y arrays multidimensionales
                                   (NCube, TPM, matriz $S$). Indexación por
                                   máscara de bits y productos matriciales vectorizados. \\
  \midrule
  \texttt{scipy}        & 1.17.0 & Clustering jerárquico
                                   (\texttt{linkage}/\texttt{fcluster}) en la variante
                                   aglomerativa \texttt{"A"} de KGeoMIP. Importado
                                   localmente para no bloquear la carga del módulo. \\
  \midrule
  \texttt{pandas}       & 2.3.3  & Lectura del Excel de pruebas via
                                   \texttt{pd.read\_excel} y escritura de resultados
                                   en CSV con \texttt{DataFrame.to\_csv}. \\
  \midrule
  \texttt{openpyxl}     & 3.1.5  & Motor \texttt{.xlsx} requerido por \texttt{pandas}.
                                   No se usa directamente en el código del proyecto. \\
  \midrule
  \texttt{pyinstrument} & 5.1.2  & Profiling estadístico; se activa con
                                   \texttt{profiler\_habilitado = True} y genera
                                   reportes de tiempo por función. \\
  \midrule
  \texttt{pyphi}        & 1.2.0  & Cálculo de la EMD-Effect (\texttt{emd\_efecto})
                                   desde \texttt{pyphi.distance}. Proporciona
                                   la norma $L_1$ entre distribuciones causales. \\
  \midrule
  \texttt{colorama}     & 0.4.6  & Coloreado de consola en \texttt{SafeLogger}:
                                   DEBUG=gris, INFO=azul, WARNING=amarillo,
                                   ERROR=rojo, CRITICAL=magenta. \\
  \midrule
  \texttt{pyttsx3}      & 2.99   & Síntesis de voz opcional en \texttt{Solution}.
                                   Por defecto desactivada en KGeoMIP y KQNodes. \\
  \bottomrule
\end{tabular}
\caption{Dependencias externas declaradas en \texttt{code/pyproject.toml}.}
\end{table}

% ─────────────────────────────────────────────────────────────────────────────
\subsection{Aspectos de Ingeniería de Software}
% ─────────────────────────────────────────────────────────────────────────────

\subsubsection*{Logging estructurado --- \texttt{SafeLogger}}

El módulo \texttt{src/middlewares/slogger.py} define la clase \texttt{SafeLogger},
un wrapper sobre \texttt{logging.Logger} con las siguientes propiedades:

\begin{itemize}
  \item \textbf{Codificación segura}: antes de emitir cada mensaje serializa
        cualquier objeto a \texttt{str} con codificación UTF-8 y sustitución
        de errores (\texttt{errors="replace"}), evitando excepciones por
        caracteres especiales en strings de partición.

  \item \textbf{Tres destinos} por instancia: (a)~archivo detallado organizado
        por fecha/hora bajo \texttt{logs/} (nivel DEBUG), (b)~archivo de última
        ejecución \texttt{last\_\{name\}.log} en la raíz de \texttt{logs/},
        y (c)~consola con colores via \texttt{colorama}.

  \item \textbf{Niveles expuestos}: \texttt{debug}, \texttt{info}, \texttt{warn},
        \texttt{error}, \texttt{critic}, \texttt{fatal}. El nivel mínimo de
        emisión está configurado a \texttt{ERROR} por defecto.

  \item \textbf{Singleton por nombre}: \texttt{logging.getLogger(name)} garantiza
        que múltiples instancias de \texttt{KGeoMIP} en el mismo proceso
        compartan el mismo logger sin duplicar handlers.
\end{itemize}

Ejemplo de uso en \texttt{KGeoMIP}:
\begin{lstlisting}
self.logger = SafeLogger("KGeoMIP")
self.logger.warn("heap agotado en _refinar_e4, k =", k)
self.logger.debug("S construida en", round(t1 - t0, 3), "s")
\end{lstlisting}

\subsubsection*{Validación de inputs}

Los parámetros se validan en las capas externas (scripts \texttt{exec\_k*.py})
antes de llegar al núcleo de cálculo:

\begin{itemize}
  \item \textbf{Máscaras binarias}: \texttt{condicion}, \texttt{alcance} y
        \texttt{mecanismo} deben ser strings de longitud $n$ con solo
        \texttt{'0'}/\texttt{'1'}. La función \texttt{\_letras\_a\_binario}
        convierte la notación del Excel (e.g., \texttt{"ABC"}) a máscara binaria
        (e.g., \texttt{"11100000"}) antes de cualquier llamada al núcleo.

  \item \textbf{Rango de \texttt{k}}: KGeoMIP acepta $1 \le k \le 5$;
        KQNodes acepta $2 \le k \le 5$. Los scripts de entrada restringen los
        valores configurables; no hay excepción explícita dentro del núcleo.

  \item \textbf{Límite $D_{\text{part}}$}: internamente se usa
        $D_{\text{part}} = \min(N, D)$ como número efectivo de dimensiones,
        evitando que $k > D$ genere bloques vacíos o índices fuera de rango.
\end{itemize}

\subsubsection*{Manejo de errores}

\begin{itemize}
  \item \textbf{\texttt{FileNotFoundError}}: si \texttt{N\{n\}\{muestra\}.csv}
        no existe en el directorio de muestras, \texttt{np.genfromtxt} lo
        propaga sin atrapar — fallo rápido visible al operador.

  \item \textbf{Heap agotado en KGeoMIP}: si el MinHeap se vacía antes de
        completar $k-1$ iteraciones, \texttt{\_refinar\_e4} emite
        \texttt{self.logger.warn} y retorna la partición parcial disponible.

  \item \textbf{Bloque de tamaño 1}: \texttt{\_qnodes\_sobre\_bloque} retorna
        \texttt{(bloque, frozenset(), 0.0)} sin invocar Queyranne, evitando
        el caso degenerado de un oracle con un único elemento.

  \item \textbf{Subsistema vacío}: si \texttt{sistema.indices\_ncubos} o
        \texttt{sistema.dims\_ncubos} está vacío tras el condicionamiento,
        \texttt{KQNodes.\_resolver} retorna \texttt{Solution} con $\varphi=0$
        y \texttt{particion="$\emptyset$"} sin entrar al refinamiento.
\end{itemize}

\subsubsection*{Profiling y debugging}

El singleton \texttt{aplicacion} centraliza la configuración global del entorno:

\begin{lstlisting}
aplicacion.profiler_habilitado = True   # activa pyinstrument
aplicacion.pagina_sample_network = "A"  # selecciona N{n}A.csv
\end{lstlisting}

Cuando el profiling está activo, los scripts \texttt{exec\_k*.py} envuelven
la llamada batch en un contexto \texttt{pyinstrument.Profiler} y emiten el
reporte de tiempos por función al finalizar. Esto permite identificar cuellos
de botella sin instrumentar el código de cálculo manualmente.

La estrategia de debugging recomendada es:
\begin{enumerate}
  \item Activar \texttt{SafeLogger} a nivel DEBUG para ver el flujo interno
        (tamaños de bloque, $\Delta\Phi$ por iteración, cache hits/misses).
  \item Usar \texttt{pyinstrument} para identificar la función dominante en tiempo.
  \item Comparar \texttt{\_calcular\_phi\_total} vs $\sum \Delta\Phi$ acumulado:
        diferencias $> 10^{-6}$ indican un bug en la lógica incremental.
\end{enumerate}

% ─────────────────────────────────────────────────────────────────────────────
\subsection{Tests Implementados}
% ─────────────────────────────────────────────────────────────────────────────

Los tests se organizan en la carpeta \texttt{tests/} de cada submódulo.

\subsubsection*{\texttt{test\_regresion\_k2\_igual\_qnodes}}

\textbf{Tipo}: regresión / integración. \\
\textbf{Objetivo}: verificar que \texttt{KQNodes(k=2)} produce exactamente el
mismo $\varphi$ que \texttt{QNodes} para todos los casos del Excel de pruebas. \\
\textbf{Casos}: 13 casos para $n \in \{5, 8, 10\}$, distintas combinaciones de
alcance y mecanismo. \\
\textbf{Criterio de éxito}: $|\Delta\varphi| < 10^{-9}$ en todos los casos. \\
\textbf{Resultado}: 13/13 verificados (ver \texttt{context/handoffs/03.md}).

\subsubsection*{\texttt{test\_regresion\_k2\_igual\_geomip}}

\textbf{Tipo}: regresión / integración. \\
\textbf{Objetivo}: verificar que \texttt{KGeoMIP(k=2)} produce exactamente el
mismo $\varphi$ que \texttt{GeoMIP} (garantía D4-01). \\
\textbf{Casos}: $n \in \{5, 8\}$ con estados iniciales \texttt{"10000"} y
\texttt{"10000000"}. \\
\textbf{Criterio de éxito}: $|\Delta\varphi| < 10^{-9}$.

\subsubsection*{\texttt{test\_ab\_c1\_vs\_c4\_n5}}

\textbf{Tipo}: A/B / calidad de heurística. \\
\textbf{Objetivo}: comparar el criterio \texttt{C4} (mín $\varphi_{\text{local}}$)
contra \texttt{C1} (máx tamaño de parte) para $k \in \{3, 4, 5\}$ con $n=5$. \\
\textbf{Métrica}: gap relativo respecto al óptimo obtenido por búsqueda exhaustiva
($S(5,k)$ pequeño). \\
\textbf{Resultado esperado}: C4 produce gap $\le$ C1 en la mayoría de los casos;
empíricamente el gap de C4 es 0 para $k=3$.

\subsubsection*{\texttt{test\_monotonicidad\_creciente}}

\textbf{Tipo}: propiedad / integración. \\
\textbf{Objetivo}: verificar que $\varphi(k+1) \ge \varphi(k)$ para todo $k$
con $n=5$, $k \in \{2, 3, 4, 5\}$. \\
\textbf{Criterio de éxito}: ningún caso viola la monotonicidad (tolerancia $10^{-9}$).

\subsubsection*{Tests de $\Delta\Phi$ incremental (D4-06)}

\textbf{Tipo}: precisión numérica / unitario. \\
\textbf{Objetivo}: verificar que \texttt{\_delta\_phi\_corte(A, B)} es igual a
$\Phi(\Pi') - \Phi(\Pi)$ calculado con \texttt{\_calcular\_phi\_total}. \\
\textbf{Criterio de éxito}: diferencia $< 10^{-6}$ para todos los cortes del
refinamiento en $n \in \{5, 8\}$.

\subsubsection*{Estrategia de validación}

La validación del proyecto opera en tres niveles complementarios:

\begin{enumerate}
  \item \textbf{Regresión $k=2$}: el caso base exacto garantiza que la extensión
        no rompe el comportamiento ya validado de GeoMIP y QNodes.
  \item \textbf{Exhaustiva para $n \le 5$}: resultados de KGeoMIP/KQNodes
        se comparan contra búsqueda exhaustiva sobre los $S(n,k)$ posibles,
        cuantificando el gap de la heurística.
  \item \textbf{Propiedades estructurales}: monotonicidad y correctitud
        del $\Delta\Phi$ incremental se verifican como invariantes del algoritmo,
        detectando regresiones en modificaciones futuras al núcleo.
\end{enumerate}
